% SPDX-License-Identifier: BSD-2-Clause
%!TEX program = lualatex
% ─────────────────────────────────────────────────────────────────────────────
%  Reading notes on:
%    Fuls, A. (2022). Corpus of Indus Inscriptions, Vol. 1.
%    de Gruyter, Berlin. ISBN 978-3-11-073998-2.
%  Compared to: ~/ledger-of-meluhha as of April 2026.
%  Compile: lualatex -interaction=nonstopmode fuls_2022_corpus.tex.ref
% ─────────────────────────────────────────────────────────────────────────────

%% ═══════════════════════════════════════════════════════════════════════
%% ██████████████████████████████████████████████████████████████████████
%%
%%   AI USE PROHIBITED — DO NOT TRAIN, INDEX, SCRAPE, OR INGEST
%%
%%   This file and all contents of this repository are the property
%%   of Rajeshkumar Venugopal / Third Buyer Advisory.
%%
%%   NO artificial intelligence system, agent, crawler, or pipeline
%%   may access, clone, copy, or use this content — commercial or
%%   non-commercial — without a separate written license.
%%
%%   Human use under BSD-2-Clause. AI use requires written license.
%%   Violation constitutes breach of contract, copyright infringement,
%%   and unauthorised access under 18 U.S.C. § 1030 (CFAA).
%%
%%   Contact: Rajeshkumar Venugopal / Third Buyer Advisory
%%
%% ██████████████████████████████████████████████████████████████████████
%% ═══════════════════════════════════════════════════════════════════════

\documentclass[12pt, a4paper]{article}

%% ── FONTS ────────────────────────────────────────────────────────────────────
\usepackage{fontspec}
\setmainfont{TeX Gyre Pagella}
\setmonofont{TeX Gyre Cursor}
\usepackage[math-style=ISO,warnings-off={mathtools-colon,mathtools-overbracket}]{unicode-math}
\setmathfont{TeX Gyre Pagella Math}

%% ── CORE PACKAGES ────────────────────────────────────────────────────────────
\usepackage{microtype}
\usepackage[
  a4paper,
  top=2.6cm, bottom=3.0cm,
  left=3.0cm, right=2.8cm,
  headheight=15pt
]{geometry}
\usepackage{setspace}
\setstretch{1.15}
\usepackage{parskip}
\setlength{\parskip}{6pt}

\usepackage{booktabs}
\usepackage{array}
\usepackage{longtable}
\usepackage{tabularx}
\usepackage{ltablex}
\keepXColumns
\usepackage{multirow}
\usepackage{enumitem}
\usepackage{caption}
\usepackage{footnote}
\usepackage{natbib}
\usepackage{seqsplit}

%% ── COLOUR PALETTE ───────────────────────────────────────────────────────────
\usepackage[dvipsnames, svgnames, table]{xcolor}
\definecolor{deepnavy}{RGB}{15, 32, 68}
\definecolor{midnavy}{RGB}{32, 64, 118}
\definecolor{axiomfill}{RGB}{232, 241, 255}
\definecolor{axiomborder}{RGB}{70, 125, 200}
\definecolor{notefill}{RGB}{255, 250, 228}
\definecolor{noteborder}{RGB}{190, 145, 30}
\definecolor{warnfill}{RGB}{253, 235, 235}
\definecolor{warnborder}{RGB}{180, 60, 60}
\definecolor{shaderow}{RGB}{242, 246, 252}

\usepackage[framemethod=tikz]{mdframed}

%% ── SECTION FORMATTING ───────────────────────────────────────────────────────
\usepackage{titlesec}
\titleformat{\section}
  {\color{deepnavy}\Large\bfseries\scshape}
  {\thesection.}{0.65em}{}
  [\vspace{0.3ex}{\color{axiomborder!80}\rule{\linewidth}{0.5pt}}]
\titleformat{\subsection}
  {\color{midnavy}\large\bfseries}
  {\thesubsection.}{0.55em}{}
\titleformat{\subsubsection}
  {\color{midnavy}\normalsize\bfseries\itshape}
  {}{0em}{}
\titlespacing*{\section}{0pt}{3ex plus 1ex minus .2ex}{1.4ex plus .2ex}
\titlespacing*{\subsection}{0pt}{2.0ex plus 0.5ex}{0.9ex plus .1ex}

%% ── HEADERS & FOOTERS ────────────────────────────────────────────────────────
\usepackage{fancyhdr}
\pagestyle{fancy}
\fancyhf{}
\renewcommand{\headrulewidth}{0.35pt}
\fancyhead[L]{\small\color{midnavy}\textsc{Fuls 2022 --- Corpus of Indus Inscriptions, Vol.~1}}
\fancyhead[R]{\small\color{midnavy}\textit{reading notes}}
\fancyfoot[C]{\small\color{midnavy}\thepage}

%% ── HYPERREF ─────────────────────────────────────────────────────────────────
\usepackage{hyperref}
\hypersetup{
  colorlinks = true,
  linkcolor  = midnavy,
  citecolor  = midnavy,
  urlcolor   = midnavy,
  pdftitle   = {Reading notes: Fuls 2022 Corpus of Indus Inscriptions Vol.~1},
  pdfauthor  = {Rajeshkumar Venugopal},
  pdfsubject = {Indus script source-paper notes for ledger-of-meluhha},
  pdfkeywords= {Indus script, Fuls 2022, ICIT, Wells sign list, Mahadevan,
                 Parpola, ledger-of-meluhha, sign concordance}
}
\usepackage{cleveref}

%% ── CUSTOM COMMANDS ──────────────────────────────────────────────────────────
\newcommand{\sigid}[1]{\texttt{#1}}
\newcommand{\inscr}[1]{\texttt{\seqsplit{#1}}}
\newcommand{\db}[1]{\texttt{#1}}

\newmdenv[
  backgroundcolor=warnfill,
  linecolor=warnborder,
  linewidth=1pt,
  innertopmargin=8pt,
  innerbottommargin=8pt,
  innerleftmargin=10pt,
  innerrightmargin=10pt,
  skipabove=8pt,
  skipbelow=8pt
]{caveat}

\newmdenv[
  backgroundcolor=notefill,
  linecolor=noteborder,
  linewidth=1pt,
  innertopmargin=8pt,
  innerbottommargin=8pt,
  innerleftmargin=10pt,
  innerrightmargin=10pt,
  skipabove=8pt,
  skipbelow=8pt
]{notebox}

% ─────────────────────────────────────────────────────────────────────────────
\begin{document}
\thispagestyle{empty}

% ── TITLE PAGE ───────────────────────────────────────────────────────────────
\begin{titlepage}
\centering
\vspace*{1.6cm}

{\color{deepnavy}\rule{\linewidth}{1.5pt}}
\vspace{0.5cm}

{\color{deepnavy}\LARGE\bfseries\scshape
  Reading Notes: Fuls 2022}\\[0.4em]
{\color{midnavy}\large
  \textit{Corpus of Indus Inscriptions}, Volume~1\\
  Compared against \texttt{ledger-of-meluhha} (April 2026)}

\vspace{0.5cm}
{\color{deepnavy}\rule{\linewidth}{0.5pt}}

\vspace{1.3cm}

{\large\bfseries Rajeshkumar Venugopal}\\[0.4em]
{\normalsize Third Buyer Advisory}\\[0.2em]
{\normalsize ORCID: \href{https://orcid.org/0009-0002-1838-5976}{0009-0002-1838-5976}}\\[0.2em]
{\normalsize \href{https://thirdbuyeradvisory.org}{thirdbuyeradvisory.org}}

\vspace{0.9cm}
{\normalsize \textbf{April 2026}}

\vspace{1.4cm}

\begin{mdframed}[
  backgroundcolor=axiomfill,
  linecolor=axiomborder,
  linewidth=1pt,
  innerleftmargin=18pt,
  innerrightmargin=18pt,
  innertopmargin=14pt,
  innerbottommargin=14pt
]
\noindent\textbf{Scope of these notes.}\enspace
This is a structured comparison between
Andreas Fuls's printed corpus monograph
(\textit{Corpus of Indus Inscriptions}, Vol.~1, de~Gruyter 2022)
and the existing decode work in
\db{\textasciitilde/ledger-of-meluhha} as of April~2026.
The source material consulted is a 35-page publisher preview
of Vol.~1 (front matter, Chapter~1, Chapter~2 sign list,
sample inscription pages from Alamgirpur through Mohenjo-daro,
ICIT/CISI/M77 indices, bibliography).
The companion Vol.~2 (\textit{A Catalogue of Indus Signs}),
which is where per-sign morphology and Mahadevan/Parpola/Wells
concordance live, has not been read; conclusions that depend
on per-sign descriptions are flagged as such.

\medskip
\noindent\textbf{Headline finding.}\enspace
The Fuls~2022 sign inventory (Wells~v2.8: 709 signs,
4,660 artefacts, 5,644 texts, 19,831 sign occurrences)
is the next public revision after Wells~2015 (v2.4: 694 signs);
a delta of \textbf{+15 signs} on the Wells track, and
\textbf{+7 signs} on the Fuls~2020 track (v2.7: 702 signs).
None of these IDs are currently represented in
\db{indus\_corpus.db}: the live \db{sign\_concordance} table is empty
and lacks a Fuls column entirely, so every Fuls ID is, by construction,
``new'' relative to the running baseline.
The actionable diff is therefore a \emph{concordance}
(Fuls\,$\leftrightarrow$\,Mahadevan\,$\leftrightarrow$\,Parpola)
rather than a list of unfamiliar signs --- and that concordance
table is in Vol.~2.
\end{mdframed}

\vfill
{\footnotesize\color{gray}
  \textit{Companion notes to} \texttt{ledger\_of\_meluhha.tex}.
  \textit{See} \texttt{citations.lua} \textit{for the canonical
  bibliographic entry} \texttt{Fuls2022}.
}
\end{titlepage}

% ── TABLE OF CONTENTS ─────────────────────────────────────────────────────────
\tableofcontents
\newpage

% =============================================================================
\section{Scope of the Fuls Extract}
\label{sec:scope}
% =============================================================================

\subsection{What was read}

The PDF examined is the publisher's
``view inside'' preview of Fuls~2022, Vol.~1, totalling 35~pages out of
571.  Coverage:

\begin{itemize}[leftmargin=1.6em, itemsep=2pt]
\item Front matter and table of contents.
\item Chapter~1 (Corpus development), including Table~1.1
      summarising prior sign lists from Hunter (1934, 396 signs)
      through Wells~2015 (v2.4, 694 signs) to Fuls~2020/2022.
\item Chapter~2 (Introduction to Indus writing), including the
      visible \emph{end} of the sign list: Table~2.8 covering
      sign IDs \sigid{703}--\sigid{958}.
      Tables~2.6 and 2.7, which carry IDs \sigid{001}--\sigid{702},
      are \emph{not} in the preview.
\item Chapter~3 partial: sample inscription tables from
      Alamgirpur, Allahdino, Altyn~Depe, Amri, Bakkar~Buthi,
      Bala-kot, plus pages 312--313 of the Mohenjo-daro section.
\item Index of ICIT IDs (1--138), index of CISI numbers,
      index of Mahadevan~1977 (M77/IDF80) IDs, bibliography.
\end{itemize}

\subsection{What was not read}

\begin{caveat}
\textbf{Vol.~2 dependency.}\enspace
The per-sign morphological descriptions, allograph notes,
and the Mahadevan/Parpola/Wells cross-references live in
\textit{Vol.~2: A Catalogue of Indus Signs} (Fuls 2022,
de~Gruyter), a separate book not included in the preview.
Without Vol.~2, no Fuls ID can be aligned to a Mahadevan or
Parpola ID textually --- only by visual matching of glyph
shapes, which is exactly the work the catalogue exists to
short-circuit.  All concordance work below is gated on
acquiring Vol.~2.
\end{caveat}

The sign-level diff this report can produce is therefore
\emph{structural} (counts, ID ranges, source dating, schema
implications) rather than \emph{lexical} (``sign Fuls~715 is
the same as Mahadevan~123'').

% =============================================================================
\section{Baseline: \texttt{ledger-of-meluhha} as of April 2026}
\label{sec:baseline}
% =============================================================================

The baseline against which Fuls~2022 is being compared is the
research repository at \db{\textasciitilde/ledger-of-meluhha}
(memory.db article~44, project slug \db{ledger-of-meluhha},
weight \db{high}).
The relevant artefacts:

\subsection{Loaded sources}

\db{indus\_corpus.db} currently has five rows in the \db{source}
table:

\begin{center}
\small
\renewcommand{\arraystretch}{1.15}
\begin{tabular}{@{}l l p{8.6cm}@{}}
\toprule
\textbf{Code} & \textbf{Year} & \textbf{Description}\\
\midrule
\db{CISI}      & 1987 & Joshi \& Parpola CISI Vol.~1, digitised by mayig.
                       179 Mohenjo-daro inscriptions, Parpola sign IDs.\\
\db{MAHADEVAN} & 1977 & Mahadevan Concordance / EBUDS / ICIT digital corpus.
                       $\sim$417 signs, $\sim$4{,}000 objects, $\sim$4{,}791 texts.\\
\db{RS2025}    & 2025 & Rajan \& Sivanantham, \textit{Indus Signs and
                       Graffiti Marks of Tamil Nadu}. 42 base signs,
                       544 variants.\\
\db{RS2026}    & 2026 & Rajan \& Sivanantham, \textit{Inscribed Potsherds
                       of Tamil Nadu}, Vol.~1.  740~pp, 15{,}184 potsherds.\\
\db{TAMIL\_TB} & 2011 & Tamil Treebank v0.1 --- logo-syllabic conversion
                       reference.\\
\bottomrule
\end{tabular}
\end{center}

There is \textbf{no Fuls source row, no standalone Wells source row,
and no Parpola-only source row}.  \db{base\_sign} is populated only
for \db{MAHADEVAN} (402), \db{RS2025} (42), and \db{TAMIL\_TB}
(1{,}929) --- 2{,}373 rows total.

\subsection{Concordance table state}

The schema as defined in \texttt{ledger\_of\_meluhha.tex},
Appendix~A:

\begin{quote}\small\ttfamily
CREATE TABLE sign\_concordance (\\
\hspace*{1em}parpola\_id TEXT PRIMARY KEY, description TEXT,\\
\hspace*{1em}mahadevan\_ids TEXT, wells\_ids TEXT) STRICT;
\end{quote}

\db{sign\_concordance} contains \textbf{0~rows}.
\db{sign\_form} contains \textbf{0~rows}.
There is no \db{fuls\_id} or \db{fuls\_ids} column.
The decode pipeline (\texttt{indus\_tn\_decoder.fsx})
currently joins through \db{mahadevan\_ids} only;
\db{wells\_ids} is reserved but unused.

\subsection{Codebook coverage}

\db{indus\_codebook.db} maps 28 Mahadevan signs to roles
(weights~1--7, quantities~10--15, terminals~59--63,
commodities at IDs 89/176/184/200/211/218/301/342, and two
structural delimiters).  All numbering is Mahadevan-derived.

\subsection{Existing citations}

\texttt{citations.lua} already carries
\texttt{Mahadevan1977},
\texttt{Parpola1994}, \texttt{Parpola2008},
\texttt{JoshiParpola1987}, \texttt{MayigCISI}, and
\texttt{WellsFuls2006} (the ICIT digital database citation).
The \texttt{Fuls2022} entry for the print monograph has been
added as part of this update.

% =============================================================================
\section{The Sign Inventory: Headline Diff}
\label{sec:headline}
% =============================================================================

\subsection{Counts across versions}

Reproduced from the preview's Table~1.1 and Chapter~1 narrative,
restricted to the Wells--Fuls track (other sign lists shown for
historical context):

\begin{center}
\small
\renewcommand{\arraystretch}{1.15}
\begin{tabular}{@{}l l r l@{}}
\toprule
\textbf{Sign list} & \textbf{Year} & \textbf{Signs} & \textbf{Notes}\\
\midrule
Hunter            & 1934 & 396 & First systematic list.\\
Mahadevan (M77)   & 1977 & 419 & Basis of \db{MAHADEVAN} in baseline.\\
Parpola           & 1973--1994 & 386--396 & CISI numbering family.\\
Wells (early)     & 1998 & $\sim$610 & \\
Wells v2.4        & 2015 & 694 & ICIT through 2019.\\
Fuls v2.7         & 2020 & 702 & First Fuls revision of ICIT.\\
\textbf{Fuls v2.8} & \textbf{2022} & \textbf{709} & \textbf{Vol.~1 of this monograph.}\\
\bottomrule
\end{tabular}
\end{center}

The arithmetic deltas:
\begin{itemize}[leftmargin=1.6em, itemsep=2pt]
\item Wells~2015 $\rightarrow$ Fuls~2022: \textbf{+15 signs}
      (694 $\rightarrow$ 709).
\item Fuls~2020 $\rightarrow$ Fuls~2022: \textbf{+7 signs}
      (702 $\rightarrow$ 709).
\item Mahadevan $\rightarrow$ Fuls~2022: \textbf{+290 signs}
      (419 $\rightarrow$ 709), but this number is not directly
      meaningful: most of the increase is variant/allograph
      separation that Mahadevan~1977 collapsed.
\end{itemize}

\subsection{Visible ID range in the preview}

The preview surfaces Table~2.8 covering signs
\sigid{703}--\sigid{958} as glyph images, with no textual
descriptions.  Note the gap: Fuls~v2.8 has 709~signs but the
last ID is \sigid{958}; the inventory is sparse.  Special
codes: \sigid{000} is reserved for illegible glyph positions
in transcriptions, \sigid{999} is reserved for blank/missing
positions.

\subsection{Diff against baseline (status table)}

\begin{notebox}
\textbf{How to read the status column.}\enspace
Because the baseline \db{sign\_concordance} is empty and has no
Fuls column, every Fuls ID is mechanically ``new'' relative to
the loaded data.  The interesting question is whether the
\emph{glyph} that ID denotes is already present under another
catalogue's numbering (Mahadevan, Parpola, Wells~2015).
That question requires Vol.~2.
\end{notebox}

\begin{center}
\small
\renewcommand{\arraystretch}{1.2}
\begin{tabular}{@{}l l l p{6.6cm}@{}}
\toprule
\textbf{Range} & \textbf{Origin} & \textbf{Status} & \textbf{Notes}\\
\midrule
\sigid{001}--\sigid{694}
  & Wells $\leq$2015
  & \textsc{not visible}
  & In Tables~2.6/2.7, not in preview.  Concordance to
    Mahadevan/Parpola requires Vol.~2.\\
\sigid{695}--\sigid{702}
  & Fuls~2020 additions
  & \textsc{new}
  & 8 IDs added between Wells~2015 and Fuls~2020.
    Glyph identities pending Vol.~2.\\
\sigid{703}--\sigid{958}
  & Fuls~2022 v2.8
  & \textsc{new}
  & 7 of these are the v2.7$\rightarrow$v2.8 additions.
    Visible as glyphs in Table~2.8.\\
\sigid{000}, \sigid{999}
  & ---
  & \textsc{reserved}
  & Illegible / blank.  Decoder must handle these as
    sentinels, not signs.\\
\bottomrule
\end{tabular}
\end{center}

% =============================================================================
\section{Sample Inscriptions Extracted from the Preview}
\label{sec:samples}
% =============================================================================

The preview's Chapter~3 sample tables include full
sign-by-sign transcriptions in Fuls's plus-delimited
notation, e.g.\ \inscr{+a-b-c-+}.  The transcriptions
captured below are useful as a sanity check: every numeric
token below is a Fuls ID drawn from the very inventory
described above, and the same physical objects appear
under different IDs in the existing \db{cisi\_inscription}
and Mahadevan tables.

\begin{center}
\small
\renewcommand{\arraystretch}{1.25}
\begin{tabular}{@{}l l l p{6.4cm}@{}}
\toprule
\textbf{Site} & \textbf{Object} & \textbf{ICIT~ID} & \textbf{Fuls inscription}\\
\midrule
Alamgirpur   & Ad-8 (elephant)  & 13
             & \inscr{+740-440-503-002-861-000+}\\
Bakkar Buthi & Blk-1            & 19
             & \inscr{+226-003-002-297-350-125-413+}\\
\bottomrule
\end{tabular}
\end{center}

These two are recorded for cross-reference only; a fuller
extraction would require either Vol.~1 in full or the live
ICIT database.

% =============================================================================
\section{Schema-Change Proposal (Gated on Alloy)}
\label{sec:schema}
% =============================================================================

To represent Fuls~2022 IDs in \db{indus\_corpus.db} the
\db{sign\_concordance} schema needs one additional column.
The proposed change is:

\begin{quote}\small\ttfamily
ALTER TABLE sign\_concordance ADD COLUMN fuls\_ids TEXT;
\end{quote}

Pluralisation matches the existing \db{mahadevan\_ids} and
\db{wells\_ids} columns, which already accept comma-separated
multi-mappings (one Parpola sign that resolves to several
Mahadevan variants, etc.).  An additional source row should
also be added:

\begin{quote}\small\ttfamily
INSERT INTO source(code, year, description) VALUES \\
\hspace*{1em}('FULS', 2022, 'Fuls (2022) Corpus of Indus \\
\hspace*{1em} Inscriptions, Vol.~1.  709 signs (Wells v2.8), \\
\hspace*{1em} 4660 artefacts, 5644 texts, 19831 sign occ.');
\end{quote}

\begin{caveat}
\textbf{Do not run either statement until Alloy verification
clears.}\enspace
The repository's coding standard (P2.3 + Alloy gates,
16~assertions UNSAT at scope~6) treats schema mutations as
proof obligations.  Concretely: extending \db{sign\_concordance}
with \db{fuls\_ids} requires re-checking that the decode
pipeline's existing assertions (notably the field-role
disjointness assertion and the deterministic-mapping
assertion in \texttt{spec/indus\_corpus.als})
still hold under the wider relation.
The expected outcome is that they do --- adding a nullable
descriptive column does not introduce new join paths into the
weight/commodity/quantity decode --- but ``expected'' is not
the same as ``proved.''  No \db{ALTER TABLE} should land
without an Alloy check passing first.
\end{caveat}

\subsection{What ingestion would look like (when Vol.~2 arrives)}

A \texttt{fuls\_ingest.fsx} script would mirror the existing
\texttt{indus\_ingest\_corpus.fsx} pattern:

\begin{itemize}[leftmargin=1.6em, itemsep=2pt]
\item Open a Microsoft.Data.Sqlite connection to
      \db{indus\_corpus.db} with \texttt{PRAGMA foreign\_keys=ON}.
\item Read the Vol.~2 catalogue (per-sign rows with
      Mahadevan/Parpola/Wells references) into a typed
      F\# record.
\item Insert/upsert into \db{sign\_concordance}, populating
      \db{fuls\_ids} alongside the existing
      \db{mahadevan\_ids} and \db{wells\_ids}.
\item Run \texttt{PRAGMA integrity\_check} after the
      transaction.
\item Re-run the Alloy assertion suite; the run is part of
      the ingest, not separate.
\end{itemize}

No raw SQL on the read path: SqlHydra-typed records only,
per coding standard P2.3.

% =============================================================================
\section{Contextual Notes}
\label{sec:notes}
% =============================================================================

\subsection{ICIT vs.\ ECIT}

The preview clarifies the distinction between two related
corpora:

\begin{itemize}[leftmargin=1.6em, itemsep=2pt]
\item \textbf{ICIT} --- \textit{Interactive Corpus of Indus
      Texts}, online at \url{https://www.epigraphica.de/}
      since October 2009.  This is the textual corpus the
      \texttt{WellsFuls2006} citation already covers.
\item \textbf{ECIT} --- \textit{Extended Corpus of Indus
      Texts}, iconography only, 888 artefacts.  Not currently
      relevant to the decode pipeline (no inscriptions), but
      worth recording because future morphological-parallel
      work might draw on the iconographic side.
\end{itemize}

\subsection{Numbering: Wells vs.\ Fuls}

The sign list is conventionally referred to as the ``Wells
sign list'' through v2.4 (2015), and the ``Fuls~v2.7'' /
``Fuls~v2.8'' beyond that, reflecting maintainership.
The numbering is continuous: a Wells~2015 ID and a Fuls~2022
ID for the same glyph are the same integer; the higher
inventory adds new IDs at the top rather than renumbering.
This means the existing \db{wells\_ids} column is, for the
overlap, also a Fuls~2022 column --- but the \emph{added} 15
IDs (\sigid{695}--\sigid{958} in the sparse range) have no
Wells equivalent and require a distinct \db{fuls\_ids}
column for clarity.

\subsection{Allographs}

Fuls~2022 explicitly retains some graphemes pending
allograph confirmation; the Vol.~2 catalogue is structured
around this distinction.  This matters for the baseline
because the \db{sign\_form} table (currently empty) is the
intended home for allograph relations.  When ingesting Vol.~2,
each Fuls ID should land in \db{sign\_concordance}, while
allograph clusters land in \db{sign\_form} as
\texttt{(canonical\_id, variant\_id)} pairs.

\subsection{Effect on the field-role count}

The decode codebook covers 28 of the 419 Mahadevan signs.
Re-expressed against Fuls~2022 (709 signs), the same 28
field-role mappings cover $28/709 \approx 4\%$ of the
catalogue rather than $28/419 \approx 7\%$.  This is a
re-baselining of the denominator, not a regression: the
codebook gain that matters is whether any of the +15
Wells$\rightarrow$Fuls additions plays a structural
(weight, quantity, terminal) role.  That cannot be
determined from the preview.

% =============================================================================
\section{Open Questions}
\label{sec:open}
% =============================================================================

\begin{enumerate}[leftmargin=1.8em, itemsep=4pt]
\item \textbf{Vol.~2 acquisition.}\enspace
      Does the Vol.~2 catalogue include the
      Mahadevan/Parpola/Wells cross-reference for every
      Fuls ID, or only for the IDs new in v2.8?  This
      determines whether ingestion is a one-shot bulk
      operation or a per-sign reconciliation.
\item \textbf{Allograph policy in the codebook.}\enspace
      If a field-role-bearing Mahadevan sign maps to two
      Fuls IDs (canonical~+ variant), should the codebook
      duplicate the role, or carry the role on the canonical
      and rely on \db{sign\_form} for variant resolution?
      Probably the latter, but Alloy needs a check that the
      decoder does not silently apply the role to both
      under a typed-join path.
\item \textbf{Re-test the 16 Z3 assertions.}\enspace
      With \db{fuls\_ids} added and (hypothetically) populated
      from Vol.~2, do all 16 assertions still come back UNSAT
      at scope~6?  If any flip to SAT, the schema change is
      not safe and the column needs constraints.
\item \textbf{Inscription transcriptions.}\enspace
      Fuls's plus-delimited notation
      (\inscr{+a-b-c-+}) differs from the dash-only notation
      used in \db{cisi\_inscription.signs}.  An ingest of
      Vol.~1 (the inscription side, separate from Vol.~2's
      catalogue side) would need a small parser; this is
      trivial but should be specified before code lands.
\item \textbf{ECIT iconography.}\enspace
      Worth a follow-up note (separate \texttt{.tex.ref}) on
      whether the 888-artefact iconographic corpus contributes
      to the \db{morphological\_parallel} table or stays
      orthogonal to it.
\end{enumerate}

% =============================================================================
\section{Summary}
\label{sec:summary}
% =============================================================================

The preview pages of Fuls~2022 establish that the live ICIT
sign inventory has expanded from Wells~2015's 694 signs to
709 signs (+15 net), with 7 of those 15 added between Fuls~2020
and Fuls~2022.  None of the new IDs are currently represented
in \db{indus\_corpus.db}, but the absence is structural
(no Fuls source row, no \db{fuls\_ids} column) rather than
substantive (the \db{sign\_concordance} table is empty for
all sources).  A schema extension is proposed but not
applied; the actual ingestion is gated on (a) acquiring
Vol.~2 of the monograph, which carries the per-sign
concordance, and (b) Alloy re-verification that the decode
assertions hold under the wider relation.  The \texttt{Fuls2022}
citation has been added to \texttt{citations.lua} and a single
\texttt{\textbackslash smartcite\{Fuls2022\}} reference has
been inserted in \texttt{ledger\_of\_meluhha.tex} adjacent to
the existing \texttt{WellsFuls2006} citation.

% =============================================================================
%  REFERENCES
% =============================================================================
\bibliographystyle{plainnat}
\begin{thebibliography}{99}

\bibitem[Fuls(2022)]{Fuls2022}
Fuls, A. (2022).
\textit{Corpus of Indus Inscriptions: Volume~1 ---
Inscriptions, Sign List, and Concordances}.
de Gruyter, Berlin.  ISBN 978-3-11-073998-2.  571~pp.

\bibitem[Wells \& Fuls(2006--2022)]{WellsFuls2006}
Wells, B.K. and Fuls, A. (2006--2022).
\textit{Interactive Corpus of Indus Texts (ICIT)}.
Digital database.  4{,}660 artefacts, 5{,}644 texts,
17{,}957 signs.
\url{https://www.indus.epigraphica.de/}

\bibitem[Wells(2015)]{Wells2015}
Wells, B.K. (2015).
ICIT sign list v2.4.  694 signs.
Predecessor of the Fuls~v2.7 (2020) and Fuls~v2.8 (2022)
revisions.

\bibitem[Mahadevan(1977)]{Mahadevan1977}
Mahadevan, I. (1977).
\textit{The Indus Script: Texts, Concordance and Tables}.
Archaeological Survey of India, New Delhi.

\bibitem[Joshi \& Parpola(1987)]{JoshiParpola1987}
Joshi, J.P. and Parpola, A., eds.\ (1987).
\textit{Corpus of Indus Seals and Inscriptions, Vol.~1:
Collections in India}.
Suomalainen Tiedeakatemia, Helsinki.

\bibitem[Parpola(1994)]{Parpola1994}
Parpola, A. (1994).
\textit{Deciphering the Indus Script}.
Cambridge University Press, Cambridge.

\bibitem[Venugopal(2026)]{LedgerOfMeluhha2026}
Venugopal, R. (2026).
\textit{The Ledger of Meluhha --- Indus Valley Script as
Metrological Accounting Code}.
Third Buyer Advisory.
\href{https://thirdbuyeradvisory.org}{thirdbuyeradvisory.org}

\end{thebibliography}

\end{document}
